% Part: incompleteness
% Chapter: introduction
% Section: historical-background

\documentclass[../../../include/open-logic-section]{subfiles}

\begin{document}

\olfileid{inc}{int}{bgr}

\olsection{پیشینهٔ تاریخی}

در این بخش به‌اختصار تحولاتی تاریخی را بررسی می‌کنیم که قضایای
ناتمامیت را در بستر خود قرار می‌دهند. به‌ویژه، مروری بسیار فشرده بر
تاریخ منطق ریاضی عرضه می‌کنیم و سپس چند کلمه دربارهٔ تاریخ مبانی
ریاضیات می‌گوییم.

\begin{digress}
  عبارت «منطق ریاضی» دوپهلو است. می‌توان واژهٔ «ریاضی» را وصف موضوع
  دانست، چنان‌که در «منطق ریاضیات»، که به اصول استدلال ریاضی اشاره
  دارد؛ یا آن را وصف روش‌ها دانست، چنان‌که در «ریاضیات منطق»، که به
  مطالعهٔ ریاضیِ اصول استدلال اشاره می‌کند. روایت پیش رو منطق ریاضی
  را در هر دو معنا، و اغلب هم‌زمان، در بر می‌گیرد.
\end{digress}

مطالعهٔ منطق اساساً با ارسطو آغاز شد که تقریباً در سال‌های
384--322~\textsc{پ.م.} می‌زیست. \emph{مقولات}، \emph{تحلیلات پیشین}
و \emph{تحلیلات پسین} او بررسی‌هایی نظام‌مند از اصول استدلال علمی،
از جمله مطالعه‌ای فراگیر و منظم دربارهٔ قیاس، در بر دارند.

منطق ارسطو در سراسر قرون وسطی بر فلسفهٔ مدرسی چیره بود؛ حتی در اواخر
سدهٔ هجدهم، کانت معتقد بود منطق ارسطو کامل و بی‌نیاز از بازنگری است.
اما نظریهٔ قیاس برای مدل‌سازی چیزی بیش از سطحی‌ترین جنبه‌های استدلال
ریاضی بسیار محدود است. یک سده پیش‌تر، لایبنیتس، هم‌عصر نیوتن،
«حسابی» کامل برای استدلال منطقی تصور کرده بود و گام‌هایی ابتدایی به
سوی طراحی آن برداشت که در اصل نسخه‌ای از منطق گزاره‌ای را توصیف
می‌کرد.

سدهٔ نوزدهم نقطهٔ عطفی برای منطق بود. جرج بول در 1854 \emph{قوانین
اندیشه} را نوشت که مطالعه‌ای جبری و فراگیر دربارهٔ منطق گزاره‌ای
داشت و چندان از ارائه‌های امروزی دور نبود. گوتلوب فرگه در 1879
\emph{مفهوم‌نگاری} (Begriffsschrift) خود را منتشر کرد که منطق
گزاره‌ای را با سورها و روابط گسترش می‌داد و بنابراین منطق مرتبه‌اول
را در بر می‌گرفت. در واقع، دستگاه‌های منطقی فرگه منطق مرتبه‌بالاتر و
مطالب بیشتری را نیز شامل می‌شدند. فرگه در \emph{قوانین بنیادی حساب}
کوشید نشان دهد همهٔ حساب را می‌توان در مفهوم‌نگاری او از فرض‌هایی
صرفاً منطقی استنتاج کرد. متأسفانه، چنان‌که راسل در 1902 نشان داد، این
فرض‌ها ناسازگار بودند. اما اگر اصل موضوعهٔ ناسازگار را کنار بگذاریم،
فرگه کم‌وبیش به‌تنهایی منطق مدرن را اختراع کرد؛ دستاوردی حیرت‌آور.
منطق سوری را متفکران جبرگرای پس از بول، از جمله پیرس و شرودر، نیز
مستقلانه پرورش دادند.

اکنون به تحولات مبانی ریاضیات می‌پردازیم. البته چون منطق در ریاضیات
نقشی مهم دارد، این تحولات با آنچه توصیف شد تعامل فراوان دارند. برای
نمونه، فرگه منطق خود را با هدف صریحِ نشان‌دادن اینکه تمام ریاضیات را
می‌توان تنها بر چارچوب منطقی او بنا کرد پرورش داد؛ به‌ویژه می‌خواست
نشان دهد ریاضیات از حقایق \emph{تحلیلیِ} پیشینی تشکیل شده است، نه
آن‌گونه که کانت معتقد بود از حقایق \emph{ترکیبیِ} پیشینی.

بسیاری زایش ریاضیات به معنای دقیق را به یونانیان نسبت می‌دهند.
\emph{اصول} اقلیدس که حدود 300 پیش از میلاد نوشته شد، با تأکیدش بر
سخت‌گیری و دقت، از پیش نمونه‌ای پخته از ریاضیات یونانی است. تعریف‌ها
و اثبات‌های \emph{اصول} اقلیدس کم‌وبیش دست‌نخورده در کتاب‌های هندسهٔ
دبیرستانی امروز باقی مانده‌اند (تا آنجا که هنوز هندسه در دبیرستان
تدریس می‌شود). این الگوی استدلال ریاضی نه‌فقط در ریاضیات، بلکه در
شاخه‌هایی از فلسفه نیز سرمشق استدلال دقیق دانسته شده است. (اسپینوزا
حتی استدلال‌های اخلاقی و دینی را به سبک اقلیدسی عرضه کرد که دیدنش
عجیب است!)

حساب دیفرانسیل و انتگرال را نیوتن و لایبنیتس در سدهٔ هفدهم ابداع
کردند. (نزاعی تند بر سر تقدم قرن‌ها ادامه داشت، اما امروزه بیشتر
پژوهشگران این دو تحول را عمدتاً مستقل می‌دانند.) این حساب، برای
نمونه، استدلال دربارهٔ مجموع‌های نامتناهیِ کمیت‌های بی‌نهایت کوچک را
در بر می‌گیرد؛ همین ویژگی‌ها انتقاد اسقف برکلی را برانگیخت که استدلال
می‌کرد باور به خدا از ریاضیات زمان او نامعقول‌تر نیست. روش‌های حساب
در سدهٔ هجدهم به‌طور گسترده به کار رفتند؛ مثلاً لئونارد اویلر با
محاسباتی شامل مجموع‌های نامتناهی به نتایجی چشمگیر رسید.

ریاضی‌دانان سدهٔ نوزدهم کوشیدند با قراردادن حساب بر بنیادی استوارتر
به انتقادهای برکلی پاسخ دهند. تلاش‌های کوشی، وایرشتراس، بولتسانو و
دیگران به تعریف‌های امروزی ما از حد، پیوستگی، مشتق و انتگرال بر حسب
«اپسیلون و دلتا»، یعنی بدون هیچ اشاره‌ای به بی‌نهایت‌کوچک‌ها،
انجامید. در اواخر آن سده، ریاضی‌دانان کوشیدند فراتر بروند و همهٔ
جنبه‌های حساب، از جمله خود اعداد حقیقی، را بر حسب اعداد طبیعی توضیح
دهند. (کرونکر: «خدا اعداد صحیح را آفرید؛ هر چیز دیگر ساختهٔ انسان
است.») ددکیند در 1872 «پیوستگی و اعداد گنگ» را نوشت و نشان داد چگونه
اعداد حقیقی را به‌عنوان مجموعه‌هایی از اعداد گویا «بسازیم» (که
می‌دانید می‌توان آن‌ها را زوج‌هایی از اعداد طبیعی دید)؛ در 1888 نیز
«اعداد چیستند و چه باید باشند؟» را نوشت که می‌کوشید اعداد طبیعی را
با اصطلاحاتی صرفاً «منطقی» توضیح دهد. کرونکر در 1887 «دربارهٔ مفهوم
عدد» را نوشت و از بازنمایی همهٔ اشیای ریاضی بر حسب اعداد صحیح سخن
گفت؛ جوزپه پئانو نیز در 1889 اصول موضوعهٔ صوری و نمادینی برای اعداد
طبیعی ارائه کرد.

پایان سدهٔ نوزدهم جسارتی تازه در برخورد با نامتناهی نیز به همراه
آورد. پیش از آن، با اشیا و ساختارهای نامتناهی (مانند مجموعهٔ اعداد
طبیعی) محتاطانه رفتار می‌شد؛ «بی‌نهایت بسیار» به معنای «هر قدر که
بخواهید» و «در حد نزدیک می‌شود» به معنای «هر قدر بخواهید نزدیک
می‌شود» فهمیده می‌شد. اما گئورگ کانتور نشان داد می‌توان نامتناهی را
به معنای ظاهری‌اش جدی گرفت. کار کانتور، ددکیند و دیگران به پیدایش
فهم عمومی و مجموعه‌نظری از ریاضیات کمک کرد که اکنون به‌طور گسترده
پذیرفته شده است.

این ما را به تحولات سدهٔ بیستم در منطق و مبانی می‌رساند. راسل در
1902 پارادوکس دستگاه منطقی فرگه را کشف کرد. زرملو در 1904 اصل
خوش‌ترتیبی کانتور را با استفاده از «اصل انتخاب» ثابت کرد؛ مشروعیت
این اصل بحث فراوانی برانگیخت. میان سال‌های 1910 و 1913 سه جلد
\emph{اصول ریاضیات} راسل و وایتهد منتشر شد و برنامهٔ فرگه برای
استوارکردن ریاضیات بر بنیاد منطقی را گسترش داد. متأسفانه راسل و
وایتهد ناچار شدند دو اصل را بپذیرند که توجیهشان به‌عنوان اصولی صرفاً
منطقی دشوار می‌نمود: اصل موضوعهٔ نامتناهی و اصل موضوعهٔ «تحویل‌پذیری».
پوانکاره در دههٔ 1900 از کاربرد «تعریف‌های غیرمحمولی» در ریاضیات
انتقاد کرد و براوئر در دههٔ 1910 پیشنهاد بازبنیادگذاری همهٔ ریاضیات
بر پایه‌ای «شهودگرایانه» را مطرح کرد که از قانون طرد شق ثالث
($!A\lor\lnot !A$) پرهیز می‌کرد.

به‌راستی روزگار عجیبی بود!{} برنامهٔ فروکاستن همهٔ ریاضیات به منطق
امروزه «منطق‌گرایی» نام دارد و عموماً به دلیل دشواری‌های گفته‌شده
شکست‌خورده تلقی می‌شود. برنامهٔ پرورش ریاضیات بر حسب ساخت‌های ذهنیِ
شهودگرایانه «شهودگرایی» نام دارد و محدودیت‌هایی بیش از اندازه سخت بر
ریاضیات روزمره تحمیل می‌کند. در حوالی آغاز سده، داوید هیلبرت، یکی از
اثرگذارترین ریاضی‌دانان همهٔ اعصار، از روش‌های انتزاعی تازهٔ کانتور و
ددکیند سخت حمایت می‌کرد: «هیچ‌کس ما را از بهشتی که کانتور برایمان
آفریده است بیرون نخواهد راند.» هم‌زمان، او به انتقادهای مبنایی از
این روش‌های تازه (که اکنون، به شکلی غریب، «کلاسیک» نامیده می‌شوند)
حساس بود. او راهی پیشنهاد کرد تا هر دو خواسته برآورده شوند:
\begin{enumerate}
\item روش‌های کلاسیک را با اصول موضوعه و قواعد صوری بازنمایی کنید؛
  پرسش‌های ریاضی را به صورت !!{formula}sیی در دستگاهی اصل‌موضوعی درآورید.
\item با روش‌های امن و «متناهی» ثابت کنید این دستگاه‌های استنتاج صوری
  سازگارند.
\end{enumerate}

کار هیلبرت برای تحقق هدف نخست بسیار پیش رفت. او در 1899 این کار را
برای هندسه در کتاب مشهور \emph{مبانی هندسه} انجام داده بود. در
سال‌های بعد، او و شماری از شاگردان و همکارانش روی دیگر حوزه‌های
ریاضیات کار کردند تا همان کاری را انجام دهند که هیلبرت برای هندسه
کرده بود. خود هیلبرت دستگاه‌هایی اصل‌موضوعی برای حساب و آنالیز ارائه
کرد. زرملو نظریهٔ مجموعه‌ها را اصل‌موضوعی کرد و فرانکل، اسکولم، فون
نویمان و دیگران آن را گسترش دادند. تا میانهٔ دههٔ 1920 دو رویکرد
مدعی عنوان اصل‌موضوعی‌سازی «تمام» ریاضیات بودند: \emph{اصول ریاضیات}
راسل و وایتهد، و آنچه نظریهٔ مجموعه‌های زرملو--فرانکل نام گرفت.

هیلبرت در 1921 پروژه‌ای پژوهشی را برای رسیدن به هدف اثبات سازگاری
این دستگاه‌ها آغاز کرد. چند تن از شاگردانش، به‌ویژه برنایس، آکرمان و
بعدها گنتسن، او را در این پروژه یاری کردند. ایدهٔ اصلی این بود که
پرسش امکان !!{derivation} یک ناسازگاری در ریاضیات را به مسئله‌ای ترکیبی
دربارهٔ دنباله‌های ممکن نمادها تبدیل کنند؛ یعنی دنباله‌های ممکن
جمله‌هایی که معیار !!{derivation} درستِ مثلاً $!A\land\lnot !A$ از اصول
موضوعهٔ دستگاهی برای حساب، آنالیز یا نظریهٔ مجموعه‌ها را برآورده
می‌کنند. اثبات ناممکن‌بودن چنین دنباله‌ای از نمادها، چون خود اثباتی
ریاضی است، در این دستگاه‌های اصل‌موضوعی صوری‌شدنی می‌بود. به بیان
دیگر، جمله‌ای چون~$\OCon$ وجود می‌داشت که مثلاً می‌گفت حساب سازگار
است. افزون بر این، این جمله باید در دستگاه‌های مورد نظر اثبات‌پذیر
می‌بود، به‌ویژه اگر اثبات آن فقط به ابزارهایی بسیار محدود و «متناهی»
نیاز داشت.

هدف دوم، یعنی اینکه دستگاه‌های اصل‌موضوعیِ ساخته‌شده هر پرسش ریاضی را
تعیین کنند، به دو شیوه دقیق‌شدنی است. یک صورت‌بندی چنین است: برای هر
جملهٔ~$!A$ در زبان دستگاهی اصل‌موضوعی برای ریاضیات، یا $!A$ یا
$\lnot !A$ از اصول موضوعه اثبات‌پذیر باشد. اگر چنین می‌بود، هیچ جمله‌ای
نبود که بر پایهٔ اصول نه اثبات و نه رد شود و هیچ پرسشی نبود که اصول
تعیین نکنند. دستگاه اصل‌موضوعیِ دارای این ویژگی \emph{کامل} نامیده
می‌شود. البته شاید هنوز برای هر جملهٔ معین، تعیین اینکه کدام گزینه
برقرار است دشوار باشد. اما در اصل باید روشی برای این کار وجود داشته
باشد. در واقع، برای دستگاه‌های اصل‌موضوعی و !!{derivation}ی مورد نظر
هیلبرت، تمامیت مستلزم وجود چنین روشی می‌بود، هرچند هیلبرت این را
نمی‌دانست. تعبیر دومِ پرسش این شرط قوی‌تر است: روشی مکانیکی و محاسباتی
وجود داشته باشد که برای جملهٔ داده‌شدهٔ~$!A$ تعیین کند آیا از اصول
موضوعه اشتقاق‌پذیر است یا نه.

گودل در 1931 دو «قضیهٔ ناتمامیت» را ثابت کرد که نشان می‌دادند این
برنامه نمی‌تواند موفق شود. هیچ دستگاه اصل‌موضوعی برای ریاضیات کامل
نیست؛ به‌ویژه، جمله‌ای که سازگاری اصول موضوعه را بیان می‌کند جمله‌ای
است که نه می‌توان اثباتش کرد و نه ردش کرد.

این ضربه‌ای مرگبار بر برنامهٔ اصلی هیلبرت وارد کرد. بااین‌حال، چنان‌که
اغلب در ریاضیات رخ می‌دهد، راه‌های تازه و هیجان‌انگیزی برای پژوهش نیز
گشود. اگر یک دستگاه صوری واحد و فراگیر برای ریاضیات وجود ندارد،
پرورش دستگاه‌هایی محدودتر و بررسی آنچه در آن‌ها اثبات‌پذیر است معنا
دارد. همچنین پرورش روش‌های اثباتی کم‌محدودیت‌تر برای نشان‌دادن سازگاری
این دستگاه‌ها و یافتن راه‌هایی برای سنجش دشواری اثبات سازگاری‌شان
معنادار است. چون گودل نشان داد (تقریباً) هر دستگاه صوری پرسش‌هایی دارد
که نمی‌تواند تعیین کند، جست‌وجوی پرسش‌های «جالبی» که دستگاه صوریِ
معینی نمی‌تواند تعیین کند و یافتن میزان قوت لازم یک دستگاه برای
تعیین آن‌ها معنا دارد. منطق‌دانان تا امروز این پرسش‌ها را در رشتهٔ
ریاضی تازه‌ای به نام نظریهٔ اثبات دنبال کرده‌اند.

\end{document}
